\documentclass[12pt]{article}
\usepackage[utf8]{inputenc}
\usepackage{amsmath,amssymb,latexsym}
\usepackage{enumerate}
\usepackage{graphicx}
\usepackage[spanish]{babel}
\usepackage{xypic}

\textheight=25cm \textwidth=18cm \topmargin=-2cm \oddsidemargin=-1cm

\newcommand{\abs}[1]{\left\vert#1\right\vert}
\newcommand{\norm}[1]{\left\Vert#1\right\Vert}

\newcommand{\Co}{\mathfrak C}
\newcommand{\diag}{\mathrm{diag}}
\newcommand{\Ext}{\mathrm{ext}}
\newcommand{\F}{\mathfrak F}
\newcommand{\h}{\mathrm h}
\newcommand{\Int}{\mathrm{int}}
\newcommand{\n}{\mathfrak{n}}
\newcommand{\N}{\mathbb N}
\newcommand{\Q}{\mathbb Q}
\newcommand{\R}{\mathbb R}
\newcommand{\US}{\mathbb{S}^2}
\newcommand{\x}{\mathrm x}
\newcommand{\xk}{(\x_k)_{k\in \N}}
\newcommand{\y}{\mathrm y}
\newcommand{\z}{\mathrm z}
\newcommand{\Z}{\mathbb Z}

\newcommand{\eps}{\varepsilon}
\newcommand{\lbil}{\left <}
\newcommand{\rbil}{\right >}
\newcommand{\To}{\longrightarrow}
\newcommand{\mTo}{\longmapsto}

\begin{document}

\pagestyle{empty}
\begin{center}
{\textsc{Cálculo Diferencial e Integral III} \\
\medskip 
\large{Ejercicios I - 3 \\
Luis David Uicab Martínez \\
(Funciones entre espacios euclidianos)}}
\end{center}
\bigskip

\begin{enumerate}
    \item Analiza la convergencia de las sucesiones $\xk$ en $\R^2$ con $\x_k=(x_k,y_k)$.
        \begin{enumerate}[(1)]
            \item $x_k=r^k, y_k=r^{1/k}$ con $0<r<1$.
            \item $x_k=(1+1/k)^k, y_k=1/\sqrt{k}$.
        \end{enumerate}
    \textbf{Solución.}
    \begin{enumerate}[(a)]
        \item Puesto que $\abs{r}<1$, sabemos que $\lim_{k\to \infty}r^k\rightarrow 0$. Puesto que además, $r>0$, se sigue que $\lim_{k\to \infty}\sqrt[k]{r}\rightarrow 1$.
            En consecuencia,
            \[
                \lim_{k\to \infty}\x_k = \left(\lim_{k\to \infty} x_k,\lim_{k\to \infty} y_k\right) = (0,1).
            \]
            \item Notemos que $\lim_{k\to \infty}(1+1/k)^k=e$. Por otra parte, tenemos que $y_k=(1/k)^{1/2}$. Así, $\lim_{k\to \infty}y_k=0$. De ello, concluimos
            \[
                \lim_{k\to \infty}\x_k = \left(\lim_{k\to \infty} x_k,\lim_{k\to \infty} y_k\right) = (e,0).
            \]
            \hfill $\blacktriangle$
    \end{enumerate}


    \item Analiza la convergencia de la sucesión $\xk$ en $\R^2$ con 
        $x_k=1+(-1)^k, y_k=1 + 1/2 + \cdots + 1/k$.
        \textbf{Solución.} Notemos que $y_k$ es la sucesión de sumas parciales de la seria armónica la cual no converge. Puesto que $y_k$ no converge, concluimos que el límite de $\x_k$ no existe. $\blacktriangle$

    \newpage
    \item Analiza la existencia del l\'imite de las siguientes funciones en los puntos dados:
        \begin{enumerate}[(1)]
            \item $\displaystyle{f(x,y) = \left( \frac{x}{x^2+y^2},\frac{-y}{x^2+y^2} \right)}$, 
                $\x_0=(0,0)$.
            \item $f(x,y) = (e^x\cos(x),e^x\sen(y))$, $\x_0=(0,0)$.
        \end{enumerate} 
    \textbf{Solución.}
    \begin{enumerate}[(a)]
        \item Procedamos por contradicción. Supongamos que existe el límite de $f(x,y)$ en $\x_0$. Es decir, existe $L \in \R^2$ tal que para toda sucesión $(\x_k)_{k\in \N}$, si $\lim_{k\to \infty}\x_k=(0,0)$, se cumple que $\lim_{k\to \infty}f(\x_k)=L$.
            Tomemos $\x_k=(1/k,0)$ la cual converge a 0. Así,
            \[
                f(\x_k) = \left(\frac{\frac{1}{k}}{\left(\frac{1}{k}\right)^2+0^2},\frac{-0}{\left(\frac{1}{k}\right)^2+0^2}\right)=(k,0),
            \]
            de lo que obtenemos una contradicción, puesto que $(k,0)$ no converge. En consecuencia, el límite de $f(x,y)$ en $\x_0$ no existe.
        \item Sea $(\x_k=(x_k,y_k))_{k\in \N}$ tal que $\lim_{k\to \infty}\x_k=(0,0)$ arbitraria. De ello, se sigue que
            \[
                \lim_{k\to \infty}x_k=0 \text{ y } \lim_{k\to \infty}y_k = 0
            \] 
            Puesto que $e^x$ y $\cos x$ son funciones continuas, afirmamos que $e^x\cos x$ también es continua.
            De este hecho obtenemos que el límite de la sucesión de las imágenes de $\x_k$ en la función ya mencionada, será la función evaluada en el límite de la sucesión $\x_k$, obteniendo así que 
            \[
                \lim_{k\to \infty}e^{x_k}\cos(x_k) = e^{\lim_{k\to \infty}x_k}\cos (\lim_{k\to \infty}x_k) = e^0\cos 0 = 1.
            \]
            Por otra parte, $e^x\sin y$ es continua, dado que $\sin y$ lo es. Así,
            \[
                \lim_{k\to \infty}e^{x_k}\sin(x_k) = e^{\lim_{k\to \infty}x_k}\sin(\lim_{k\to \infty}y_k) = e^0\sin 0 = 0.
            \]
            De las igualdades anteriores, se sigue que $\lim_{k\to \infty}f(\x_k)=(1,0)$. Puesto que lo anterior se cumplió para una sucesión que tiende a 0 arbitraria, afirmamos que se cumple para toda sucesión $\x_k$ que converge a $(0,0)$, de lo que concluimos, el límite de $f(x,y)$
            en $(0,0)$ existe y además es $(1,0)$. \hfill $\blacktriangle$
    \end{enumerate}

\newpage
\item Analiza la existencia del l\'imite de la función
$\displaystyle{f(x,y) = \left( \frac{\sen(x)}{x},y\sen(1/y) \right)}$, en el punto $(0,0)$.
\textbf{Solución.} Por L'Hôpital tenemos que
\[
\lim_{x\to 0}\frac{\sin x}{x} =\lim_{x\to 0} \frac{\cos x}{1} = 1.
\]
Por otra parte, tenemos que
\begin{align*}
    -1\leq \sin\left(\frac{1}{y}\right) \leq 1 &\implies -\abs{y}\leq \sin\left(\frac{1}{y}\right) \leq \abs{y} \\
    &\implies 0 \leq \abs{\sin\left(\frac{1}{y}\right)} \leq \abs{y},
\end{align*}
luego, puesto que $\lim_{y\to 0}0=\lim_{y\to 0}\abs{y}$, por el lema del sandwich y de la desigualdad anterior, $\lim_{y\to 0}y\sin\left(\frac{1}{y}\right)=0$. Finalmente, afirmamos que
\[
    \lim_{(x,y)\to 0}f(x,y) = \left(\lim_{(x,y)\to 0}\frac{\sin x}{x},\lim_{(x,y)\to 0} y\sin\left(1/y\right)\right) = \left(1,0\right).
\]
\hfill $\blacktriangle$


\item Analiza la existencia del l\'imite de las siguientes funciones en los puntos dados:
\begin{enumerate}[(1)]
\item $\displaystyle{f(x,y) = \frac{x^3y^2}{x^6+y^4}}$, $\x_0=(0,0)$.
\item $\displaystyle{f(x,y) = \frac{x^2(y-1)}{x^4+(y-1)^2}}$, $\x_0=(0,1)$.
\end{enumerate}
\textbf{Solución.}
\begin{enumerate}[(a)]
    \item Sea $(a_k)_{k\in \N}=(x_k,x_k^{3/2})_{k\in\N}$ tal que $(a_n)_{x_k}$ con $x_k=1/k$. Así, $(a_k)_k\in \N$ converge a $(0,0)$. Notemos que
    \[
    f(a_k)=\frac{x_k^3x_k^3}{x_k^6x^6}=\frac{1}{2},
    \]
    de lo que sigue, $\lim_{k\to \infty}f(a_k)=1/2$. Por otra parte, si tomamos la sucesión $(1/k,1/k)_{k\in \N}$, la cual es convergen a 0, obtenemos que
    \[
    f\left(\frac{1}{k},\frac{1}{k}\right) = \frac{\frac{1}{k^5}}{\frac{1}{k^6}+\frac{1}{k^4}} = \frac{k}{1+k^2}.
    \]
    Así,
    \[
    \lim_{k\to \infty}f\left(\frac{1}{k},\frac{1}{k}\right)= \lim_{k\to \infty}\frac{k}{1+k^2} = 0,
    \]
    y puesto que tenemos dos sucesiones que convergen a distintos valores, afirmamos que $\lim_{\x\to \x_0}f(x,y)$ no existe.
    \item Sea $(a_k)_{k\in \N}=(1/k,1+1/k)_{n\in \N}$ donde afirmamos que $(a_k)_{k\in \N}$ converge a $(0,1)$. Luego, notemos que
    \[
    \lim_{k\to \infty}f(a_k) = \lim_{k\to \infty}\frac{1/k^3}{1/k^4+1/k^2} = \lim_{k\to \infty}\frac{k}{1+k^2} = 0.
    \]
    Ahora, notemos que $(1/k^{1/2},1+1/k)_{k\in \N}$ también converge a $(0,1)$. Luego, tenemos que
    \[
    \lim_{k\to \infty}f\left(1/k^{1/2},1+1/k\right) = \lim_{k\to \infty}\frac{1/k^2}{1/k^2+1/k^2} = \lim_{k\to \infty}\frac{1/k^2}{2/k^2} = 1/2.
    \]
    En consecuencia, dado que distan los límites, concluimos que $\lim_{\x \to \x_0}f(\x)$ no existe. \hfill $\blacktriangle$ 
    
\end{enumerate}



\newpage
\item Analiza la existencia del l\'imite de las siguientes funciones en los puntos dados:
\begin{enumerate}[(1)]
\item $\displaystyle{f(x,y) = \left( \frac{x^3-4}{x^2+1}, \frac{y^2-4}{3y-6} \right)}$, 
$\x_0=(2,2)$.
\item $\displaystyle{f(x,y) = \left( x^{3/2}, 1 - \frac{1}{2}y^2 \right)}$, $\x_0=(0,0)$.
\end{enumerate} 
\textbf{Solución.}
\begin{enumerate}[(a)]
    \item Puesto que el cociente entre polinomios es continuo, y dado que la expresión de la coordenada en $x$ está bien definida para $x=2$, afirmamos que
    \[
    \lim_{x\to 2}\frac{x^3-4}{x^2+1} = \frac{8-4}{4+1} = \frac{4}{5}. 
    \]
    Por otra parte, notemos que
    \[
    \frac{y^2-4}{3y-6} = \frac{(y+2)(y-2)}{3(y-2)} = \frac{y+2}{3}, \quad \forall y \in \R\setminus\{2\}.
    \]
    Así,
    \[
    \lim_{y\to 2}\frac{y^2-4}{3y-6} = \lim_{y\to 2}\frac{y+2}{3} = \frac{4}{3}.
    \]
    En consecuencia, de los resultados de ambos límites concluimos que
    \[
    \lim_{x\to x_0}f(x,y) = \left(\frac{4}{5},\frac{4}{3}\right).
    \]
    \item Tomemos una sucesión arbitaria $(x_k,y_k)_{k\in \N}$ contenida en el dominio de $f$ que converja a $(0,0)$. Puesto que, tanto $x^{3/2}$ y $1-y^2/2$ son funciones continuas,
    se sigue que
    \[
    \lim_{k\to \infty}x_k^{3/2}=\left(\lim_{k\to \infty}x_k\right)^{3/2} = 0,
    \]
    y
    \[
    \lim_{k\to \infty}1-y_k^2/2 = 1-\left(\lim_{k\to \infty}y_k\right)^2/2 = 1.
    \]
    Así, concluimos que $\lim_{\x\to \x_0}f(\x)=(0,1)$.
    \hfill $\blacktriangle$
    
\end{enumerate}

\newpage
\item Analiza la existencia del l\'imite de las siguientes funciones en los puntos dados:
\begin{enumerate}[(1)]
\item $\displaystyle{f(x,y,z) = \frac{xy^2}{\sqrt{x^4+y^4+z^2}+z^4}}$, $\x_0=(0,0,0)$.
\item $\displaystyle{f(x,y) = \left( \frac{\cos(x)-1}{x}, \frac{\sen(y)}{y} \right)}$, 
$\x_0=(0,0)$.
\end{enumerate}
\textbf{Solución.}
\begin{enumerate}[(a)]
    \item En principio, notemos que
    \begin{align*}
        \sqrt{x^4+y^4+z^2} + z^4  &\geq \sqrt{x^4+y^4}. \tag{1}
    \end{align*}
    Por la desigualdad de la media aritmética-media geométrica, tenemos
    \begin{align*}
        \frac{x^4+y^4}{2}\geq \sqrt{x^4y^4} &\implies x^4+y^4 \geq 2x^2y^2 \\
        &\implies \sqrt{x^4+y^4} \geq \sqrt{2}\abs{x}\abs{y}.
    \end{align*}
    Así, por la desigualdad anterior y la mencionada en $(1)$ se sigue que
    \[
        \sqrt{x^4+y^4+z^2} + z^4 \geq \sqrt{2}\abs{x}\abs{y},
    \]
    y en consecuencia,
    \[
    0\leq \abs{f(x,y,z)} = \frac{\abs{x}y^2}{\sqrt{x^4+y^4+z^2} + z^4} \leq \frac{\abs{x}y^2}{\sqrt{2}\abs{x}\abs{y}} = \frac{\abs{y}}{\sqrt{2}}. \tag{2}
    \]
    Notemos que la última expresión está particularmente bien definida si y solo si $x,y\neq 0$. Sea $(x_k,y_k,z_k)_{k\in \N}\subset \text{dom}(f)$ convergente a 0.
    Si $x_k=0$ o $y_k=0$, se tiene que
    \[
    f(0,y_k,z_k)= \frac{0\cdot y_k^2}{\sqrt{x_k^4+y_k^4+z_k^2} + z_k^4} = 0 = \frac{x_k\cdot 0^2}{\sqrt{x_k^4+y_k^4+z_k^2} + z_k^4} = f(x_k,0,z_k),
    \]
    y en consecuencia, para tales casos, $\lim_{k\to \infty}f(x_k,y_k,z_k)=(0,0,0)$. Procederemos a enfocarnos en las sucesiones convergentes a 0 donde $x_k,y_k\neq 0$.
    Para dichas sucesiones, la última desigualdad en $(2)$ está bien definida. Notemos que, puesto que $(x_k,y_k,z_k)\rightarrow (0,0,0)$, en particular $y_k \rightarrow 0$.
    De ello, $\lim_{k\to \infty}\abs{y_k}/\sqrt{2}=0$. Luego, del límite ya mencionado, del lema del sandwich y de la desigualda en $(2)$
    \[
    \lim_{k\to \infty}\abs{f(x_k,y_k,z_k)}=(0,0,0),
    \]
    que por el ejercicio número 8, es equivalente a decir
    \[
    \lim_{k\to \infty}f(x_k,y_k,z_k)=(0,0,0).
    \]
    De la ya dicho, concluimos que $\lim_{\x\to \x_0}f(\x)=0$. \hfill $\blacktriangle$


    \newpage
    \item Notemos que por L'Hôpital 
    \[
    \lim_{x\to 0}\frac{\cos(x)-1}{x} = \lim_{x\to 0}\frac{-\sin(x)}{1} = 0,
    \]
    y
    \[
    \lim_{y\to 0}\frac{\sin(y)}{y} = \lim_{y\to 0}\frac{\cos(y)}{1}=1,
    \]
    y en consecuencia, $\lim_{\x\to \x_0}f(\x)=(0,1)$. \hfill $\blacktriangle$
\end{enumerate}

\item Supongamos que $\Omega\subset \R^n$, $f\in \F(\Omega,\R^m)$ y 
$\x_0\in \Omega'$. Prueba que 
$\displaystyle{\lim_{\x\to \x_0} f(\x) = 0}$ si y s\'olo si 
$\displaystyle{\lim_{\x\to \x_0} \norm{f(\x)} = 0}$. \\
\textbf{Demostración.} Sea $\varepsilon > 0$ arbitrario. Puesto que $\lim_{\x\to \x_0}f(\x) = 0$ se sigue que
existe $\delta > 0$ tal que
\[
\norm{\x-\x_0}<\delta \implies \norm{f(\x)}< \varepsilon.
\]
Como $\norm{\x}=\abs{\norm{\x}}$ y $\norm{f(\x)}=\abs{\norm{f(\x)}}$, es inmediato que
\[
\norm{\x}<\delta \implies \norm{f(\x)}< \varepsilon \iff \abs{\norm{\x}}<\delta \implies \abs{\norm{f(\x)}}< \varepsilon. \tag{1}
\]
y puesto que la segunda implicación es la definición de que $\lim_{\x\to \x_0}\abs{f(x)}=0$, de la equivalencia en $(1)$ se sigue el resultado. \hfill $\blacktriangle$

\newpage
\item Supongamos que $f:\Omega\subset \R^n\To \R^m$, $\x_0\in \R^n$ y 
$L\in \R^m$. Definamos $\Omega_0=\Omega-\x_0:=\{ \x-\x_0 : \x\in \Omega \}$ y $g:\Omega_0\To \R^m$ por $g(\h)=f(\x_0+\h)$, para cada $\h\in \Omega_0$. Demuestra que:
\begin{enumerate}[(1)]
\item $\x_0\in \Omega'$ si y s\'olo si $0\in \Omega'_0$.
\item $\displaystyle{L=\lim_{\x\to \x_0} f(\x)}$ si y s\'olo si 
$\displaystyle{L=\lim_{\h\to 0} g(\h)}$.
\end{enumerate}
\textbf{Demostración.}
\begin{enumerate}[(a)]
    \item Notemos que
    \begin{align*}
        0 \in \Omega_0 &\iff 0=\x_0-\x_0 \in \Omega_0 \\
        &\iff \x_0 \in \Omega.
    \end{align*}
    \item Sea $\varepsilon>0$. Supongamos que $\lim_{\x\to\x_0}f(\x)$. Así, existe $\delta >0$ tal que, para $\x\in \Omega$:
    \[
        \norm{\x-\x_0}<\delta \implies \norm{f(\x)-L}<\varepsilon. \tag{1}
    \]
    Notemos que si tomamos $h=\x-\x_0$, de la implicación anterior, se sigue que
    \begin{align*}
        \norm{h}=\norm{\x-\x_0}<\delta \implies \norm{g(h)-L} &=\norm{f(x_0+\h)-L}\\
        &=\norm{f(\x_0+\x-\x_0)-L} \\
        &=\norm{f(\x)-L}<\varepsilon,
    \end{align*}
    de lo que afirmamos
    \[
    \lim_{\h\to 0}g(\h)=L.
    \]
    Tomemos ahora como hipótesis la igualdad anterior. Así, existe $\delta > 0$ tal que, para $\h \in \Omega_0$:
    \[
        \norm{h}<\delta \implies \norm{g(h)-L}<\varepsilon.
    \]
    Puesto que $h \in \Omega_0$, se sigue que $h=\x-\x_0$ para algún $\x \in \Omega$. Así, tenemos que
    \begin{align*}
        \norm{\x-\x_0}=\norm{h}<\delta \implies \norm{f(\x)-L} &= \norm{f(\x_0+\x-\x_0)-L} \\
        &= \norm{f(\x_0+h)-L} \\
        &= \norm{g(h)-L} < \varepsilon.
    \end{align*}
    En consecuencia, concluimos que
    \[
    \lim_{\x\to\x_0}f(\x)=L,
    \]
    con lo que damos el resultado por demostrado. \hfill $\blacktriangle$

\end{enumerate}

\newpage
\item Si $L:\R^n\To \R^m$ es una funci\'on lineal, prueba que si
\[ \lim_{\x\to 0} \frac{\norm{L(\x)}}{\norm{\x}} = 0 \]
entonces $L$ es la funci\'on constante cero.
\begin{enumerate}
    \item Sea $\{e_1,...,e_n\}$ la base canónica de $\R^n$. Así, para $\x=(x_1,...,x_n) \in \R^n$, se sigue que
    \[
    L(\x) = L\left(\sum_{i=1}^{n}x_ie_i\right) = x_iL(e_i), \tag{1}
    \]
    esto último por linealidad de $L$. Sea $k$ tal que
    \[
    \norm{L(e_k)} = \max\{\norm{L(e_i)}:i\in \{1,...,n\}\}.
    \]
    Luego, puesto que $\lim_{\x\to 0}\norm{L(\x)}/\norm{\x}= 0$, se sigue que, dado $\varepsilon > 0$ arbitrario,
    existe $\delta>0$ tal que
    \[
    \norm{\x}<\delta \implies \frac{\norm{L\left(\norm{x}\right)}}{\norm{\x}} < \varepsilon.
    \]
    Puesto que $\norm{\delta e_k/2} < \delta$, de la implicación previa, de la linealidad de $L$ y del hecho de que la
    norma saca escalares, obtenemos que
    \begin{align*}
        \norm{L(e_k)} &= \frac{\norm{L(e_k)}}{\norm{e_k}} \\
        &= \frac{\delta/2}{\delta/2}\cdot\frac{\norm{L(e_k)}}{\norm{e_k}} \\
        &= \frac{\norm{\delta L(e_k)/2}}{\norm{\delta e_k/2}} \\
        &= \frac{\norm{L(\delta e_k/2)}}{\norm{\delta e_k/2}} \\
        &< \varepsilon.
    \end{align*}
    Puesto que la igualdad anterior se cumple para todo $\varepsilon > 0$, afirmamos que $\norm{L(e_k)}=0$. Luego,
    como $\norm{L(E_k)}\geq \norm{L(e_i)} \geq 0$ para toda $i \in \{1,...,n\}$, se sigue
    \[
    \norm{L(e_i)}=0, \quad \forall i \in \{1,...,n\}.
    \]
    Lo anterior pasa si y solo si $L(e_i)=0$, con lo que, por la igualdad en $(1)$ se sigue el resultado. \hfill $\blacksquare$

\end{enumerate}
\end{enumerate}

\end{document}
